% 第5章图：SVD求解的坐标变换示意，m=4、d=3、r=2。符号示意，无数值。
% b -> c=U^T b（前r个可拟合，其余留在残差）-> z_i=c_i/sigma_i（i<=r），零奇异方向 z_3=0 -> x=Vz=V_r z_{1:r}。
\begin{tikzpicture}[x=1mm,y=1mm,>=Stealth,
  cell/.style={rectangle,rounded corners=1.8pt,draw=FigureStroke,fill=white,line width=.7pt,minimum width=13mm,minimum height=6mm,inner sep=1pt,font=\figurefont\fontsize{8}{9.6}\selectfont},
  fit/.style={cell,draw=FigureBlue,fill=FigureBlueFill},
  res/.style={cell,draw=FigureMuted,dashed,fill=white,text=FigureMuted},
  zero/.style={cell,draw=FigureGreen,fill=FigureGreenFill},
  head/.style={font=\figurefont\fontsize{8}{9.6}\selectfont},
  note/.style={font=\FigureNoteFont\fontsize{7.5}{9}\selectfont,text=FigureMuted}]
\def\dx{27}\def\dy{7}
% column 1: b
\foreach \i in {1,...,4}{\node[cell] (b\i) at (0,{-(\i-1)*\dy}) {$b_{\i}$};}
% column 2: c = U^T b
\foreach \i in {1,2}{\node[fit] (c\i) at (\dx,{-(\i-1)*\dy}) {$c_{\i}$};}
\foreach \i in {3,4}{\node[res] (c\i) at (\dx,{-(\i-1)*\dy}) {$c_{\i}$};}
% column 3: z
\node[fit] (z1) at ({2*\dx},0) {$c_1/\sigma_1$};
\node[fit] (z2) at ({2*\dx},{-\dy}) {$c_2/\sigma_2$};
\node[zero] (z3) at ({2*\dx},{-2*\dy}) {$0$};
% column 4: x
\foreach \i in {1,2,3}{\node[cell] (x\i) at ({3*\dx},{-(\i-1)*\dy}) {$x_{\star,\i}$};}
% headings
\node[head] at (0,7.5) {$\bm b$};
\node[head] at (\dx,7.5) {$\bm U^\top\bm b$};
\node[head] at ({2*\dx},7.5) {$\bm z$};
\node[head] at ({3*\dx},7.5) {$\bm x_\star=\bm V\bm z$};
% arrows
\foreach \a/\b/\l in {b/c/{$\bm U^\top$},c/z/{$\div\sigma_i$},z/x/{$\bm V$}}{
  \draw[->,FigureStroke,line width=.8pt] ([xshift=1.2mm]\a2.east|-0,{-0.5*\dy})--++({\dx-15.4},0)
    node[midway,above,font=\figurefont\fontsize{8}{9.6}\selectfont] {\l};
}
% notes
\node[note,anchor=north] at (\dx,{-3*\dy-3.6}) {留在残差中};
\node[note,anchor=north] at ({2*\dx},{-2*\dy-3.6}) {零奇异方向置零};
\end{tikzpicture}
